\section{Directions}
\label{sec:directions}

This section maps what the comparison opens up. Items are grouped by how much
has to be invented before they can be built, and each carries the same four
labels used in Section~\ref{sec:learn}: what it is, what it buys, what it
costs, and which invariant it puts at risk.

A caution first, which is this project's own standing discipline and is worth
inheriting: \emph{nothing here is scheduled by virtue of being written down}.
Each item is recorded with the trigger that would signal it has become
necessary, so that a future reader recognises the moment instead of
pre-building for it.

\subsection{The organising picture}

Almost everything below is a point on one line, so it is worth drawing the
line first.

\begin{figure}[h]
\centering
\begin{tikzpicture}[font=\small]
% the spectrum bar
\draw[very thick, -{Stealth[length=3mm]}] (0,0) -- (13.2,0);
\foreach \x/\lab/\sub in {
   0.6/{asserted DAG}/{nothing materialised},
   4.6/{+ selected meets}/{MDL, or workload},
   9.6/{full lattice}/{every closed set}}
  {\filldraw[odblue] (\x,0) circle (2.4pt);
   \node[anchor=south, align=center, font=\small\bfseries]
     at (\x,0.18) {\lab};
   \node[anchor=north, align=center, font=\scriptsize\itshape, text=odgrey]
     at (\x,-0.18) {\sub};}
\node[anchor=north, align=center, font=\scriptsize, text=odrust]
  at (0.6,-0.95) {\odag{} today\\cheap, canonical\\query pays the rest};
\node[anchor=north, align=center, font=\scriptsize, text=odgreen]
  at (4.6,-0.95) {the interesting middle\\directions D7 and D8};
\node[anchor=north, align=center, font=\scriptsize, text=odrust]
  at (9.6,-0.95) {\fca{}\\complete, exponential\\lookup pays nothing};
\draw[decorate, decoration={brace, amplitude=5pt}, odgrey]
  (1.0,0.95) -- (9.2,0.95)
  node[midway, above=5pt, font=\scriptsize\itshape, text=odgrey]
  {one family, parameterised by which meets are materialised};
\end{tikzpicture}
\caption{\odag{} and \fca{} are the two endpoints of a single family, not two
rival data structures. A member of the family is: the asserted graph, plus a
set of materialised meets, plus a storage mode per cone. \odag{} materialises
none and pays at query time; \fca{} materialises all and pays in space. The
research question is which middle point a given workload wants.}
\label{fig:spectrum}
\end{figure}

Seeing it this way reframes several apparently separate questions as one.
``Which concepts should MDL learn?'', ``which cones should be indexed?'',
``which meets should the query planner materialise?'' and ``which categories
should a user be prompted to create?'' are the same question asked by four
different subsystems.

\subsection{Buildable now}

These need no new theory and threaten no invariant.

\paragraph{D1. Re-describe dimensions as an interval pattern structure.}
\emph{Cost: an afternoon of writing.} Section~\ref{sec:learn-patterns} showed
\odag{}'s typed values are an instance of \citet{ganter2001}. Recording that
in the design documents converts an engineering hunch into a known-correct
construction, gives the multi-coordinate question a ready answer (product
semilattice, componentwise meet), and supplies the projection theory that the
future per-dimension index will need. Highest value-to-cost ratio in the
paper.

\paragraph{D2. An ordinal kind.} \emph{Trigger: someone files ranks.} The
scale table of Section~\ref{sec:learn-scaling} has a hole where
\textsf{beginner} $<$ \textsf{intermediate} $<$ \textsf{expert} should be.
Declaring an ordinal chain is monotone and canonicalisable, so it passes both
admissibility axes. Build it with the arithmetic deliberately \emph{absent},
because the standing hazard with ordinal data is people averaging ranks
\citep{stevens1946}.

\paragraph{D3. Reduced labelling in the visualiser.} \emph{Cost: a rendering
change.} Compare Figures~\ref{fig:lattice} and~\ref{fig:ontodag}: the same
information, far less ink. Worth most on stores with many typed values.

\paragraph{D4. An implication lint.} \emph{Trigger: users mis-filing things
that the store's own contents would have caught.} Compute extensional
implications over a pinned root into a derived store; surface them as prompts.
Never store them as claims (\S\ref{sec:learn-implications}). The natural first
version is single-premise implications only, which are cheap --- one closure
per category.

\paragraph{D5. Cone bitmaps behind \cmd{get}.} \emph{Trigger: queries
measurably hot.} An in-memory bitmap per category, intersected with a bitwise
\textsc{and}. Bounded, dependency-free (Python integers suffice), invisible in
the schema, and it has a free test oracle: the population count of a cone
bitmap must equal the maintained \verb|descendant_count|. This is the first
step of the semantic-codes programme and the one that needs no decisions. When
it outgrows plain integers the prior art is well mapped: compressed bitmaps
\citep{lemire2018} for the representation, and interval labelling
\citep{agrawal1989} for the trick of numbering nodes so that a cone becomes a
few contiguous ranges.

\subsection{Medium term: needs a design, not a discovery}

\paragraph{D6. Attribute exploration with a language model as the expert.}
\label{sec:dir-exploration}
\emph{Trigger: available now; the blocker is willingness, not knowledge.}

This is the most promising item in the paper, and
Sections~\ref{sec:learn-exploration} and~\ref{sec:agents-review} argued why:
it is the only proposal here that makes review of machine-written knowledge
\emph{bounded, ordered and terminating}. A concrete v1 for the single-writer
case:

\begin{enumerate}[leftmargin=1.7em,itemsep=1pt]
\item Pin a root. Everything that follows is relative to it, so the session is
  reproducible and citable.
\item View the store as the context $\ctx(P)$ of Section~\ref{sec:dictionary}
  and compute the next unconfirmed implication by NextClosure
  \citep{ganter1984}.
\item Ask the agent, over the MCP surface, in domain language: \emph{is every
  X also a Y?}
\item \textbf{Yes} $\to$ record a signed confirmation in the provenance store.
  If the implication has a single-category premise it is also a subsumption,
  so offer it as an ordinary \cmd{put}.
\item \textbf{No} $\to$ require a counterexample, and file it with \cmd{put}.
  The store grows by exactly the object that refutes the question.
\item Terminate when no unconfirmed implications remain. Present the human
  with the confirmed base --- which is minimal and non-redundant by
  construction.
\end{enumerate}

Two properties make this more than a demo. The counterexample requirement
means a \emph{No} is as informative as a \emph{Yes} --- the store improves
either way, which is not true of ordinary agent review. And because every
answer is a signed claim against a pinned root, a later reader can see
precisely which questions this author answered, and re-run the exploration
under a different author to compare.

Open sub-problem: a wrong expert. Classical exploration assumes infallibility;
attribution degrades the guarantee honestly to ``complete relative to what
this author asserted'', which is the right weakening. A second agent
re-answering the same question sequence, with disagreements surfaced, is the
obvious cheap defence and would be worth measuring.

\paragraph{D7. Workload-optimal materialisation.}
\label{sec:dir-materialisation}
\emph{Trigger: a real query log.} Figure~\ref{fig:spectrum}'s middle, stated
as an optimisation problem: choose a set $M$ of materialised meets and a
storage mode per cone to minimise
\[
  \mathbb{E}_{q \sim Q}\bigl[\text{retrieval cost}(q; M)\bigr]
  \;+\; \lambda_{\text{space}}\cdot \text{storage}(M)
  \;+\; \lambda_{\text{maint}}\cdot \text{update cost}(M).
\]
This is a known problem three times over: view-materialisation selection on a
lattice \citep{harinarayan1996}, whose greedy solution is near-optimal by
submodularity; retrieval-aware MDL, which is the same objective seen from the
learner's side; and adaptive indexing \citep{idreos2007}, which supplies the
online algorithm --- every query that computes a residual intersection has
just computed a candidate meet, so admit it when expected reuse covers its
cost and evict when cold.

Two \odag{}-specific notes. The selectivity statistics the greedy algorithm
needs already exist and are exact (\verb|descendant_count|). And the space
term is not a tuning constant on a rent-charging network: it is postage
(\S\ref{sec:crypto-economics}).

\emph{The soundness condition, which is easy to miss:} materialised meets may
be used to \emph{prune} candidates but not to \emph{generate} them, unless
placement is made canonical --- because of the meet-substitution trap
(\S\ref{sec:meet-trap}). A node under $A$ and $B$ is not the meet of $A$ and
$B$.

\paragraph{D8. MDL-learned concepts, with a promotion path.}
\label{sec:dir-mdl}
\emph{Trigger: a store big enough that nobody can see its structure.} Run the
learner \citep{mdlfca} over a pinned root; keep its output in a derived store;
let a human or an accountable agent promote a learned category into the shared
graph as an ordinary claim recording the root it was derived from. The
layering is what keeps this safe (\S\ref{sec:learn-mdl}), and the promotion
step is what keeps it attributable.

The interesting scientific question underneath is whether MDL-selected
concepts agree with the ones people actually name. That is a measurable
question --- take a hand-built ontology, hide the intermediate categories,
learn from the leaf assignments, compare --- and nobody appears to have run
it at scale.

\paragraph{D9. Merkle-ised cone commitments.} \emph{Trigger: somebody needs a
subsumption verified on chain.} A per-node cone-commitment index, in its own
store, manifest-pinned, deterministic from the canonical form. It buys
path-sized order-free positive proofs, per-category agreement by one hash
comparison, and the one proof shape a contract can check
(\S\ref{sec:crypto}). It can never be the base representation --- names are
identity, and a Merkle address changes whenever anything below the node
changes --- so it is an index over the named graph, forever.

\paragraph{D10. The thin browser client.} \emph{Trigger: a published store
anyone should be able to query without downloading.} Most pieces exist: a
lazy reader that fetches records as a query walks them, published cone
summaries that collapse broad queries to a handful of fetches, and a base
install that is pure Python. The remaining work is batching the fetch frontier
and bridging to network access from inside a browser sandbox.

\subsection{Longer term: genuinely open}

\paragraph{D11. Concurrent exploration.} Two agents exploring two replicas ask
different questions and confirm different things. Since exploration assumes a
fixed universe and \odag{} assumes a merging one, reconciling them is open.
A promising angle: exploration confirmations are claims, claims merge
conflict-free, and the \emph{question} sequence is a derived, local artefact
--- so perhaps only the answers need to converge, and each replica may
legitimately have asked a different route to them.

\paragraph{D12. Relations.} The biggest wall. ``Alice authored Doc1'' is not
subsumption, and no amount of cleverness makes it one. Two things are known.
The \fca{}-native answer, Relational Concept Analysis
\citep{rouanehacene2013}, computes a global fixpoint, which conflicts with
local canonical form. And the description-logic answer \citep{baader2003}, an EL-shaped
extension, is \emph{monotone} --- so merge would survive --- but makes
subsumption a global inference, so canonicalising up to logical equivalence
means canonicalising an entailment closure. That is the axis that breaks, and
naming it precisely is the contribution of Section~\ref{sec:two-canonical-forms}:
the guarantee would silently degrade from ``equal knowledge, equal root'' to
``equal syntax, equal root'', which is the weaker thing consumers would
mistakenly keep relying on. Anyone attempting this should treat cheap semantic
canonicalisation of the fragment as the \emph{primary} research problem, not
tractability of reasoning.

There is a precise precedent for drawing a line this way: EL++ retains
polynomial subsumption only under ``p-admissible'' concrete domains
\citep{baader2005} --- the description-logic community's version of \odag{}'s
rule that only exact-arithmetic kinds may enter the canonical order.

\paragraph{D13. Zero-knowledge subsumption.} \emph{Trigger: a
privacy-demanding counterparty.} Positioned in Section~\ref{sec:crypto-zk};
the encouraging part is that the properties a circuit wants --- determinism,
one primitive, no floats --- were already bought for other reasons.

\paragraph{D14. Products of DAGs, and factorial codes.} This one returns to
the origin. Section~\ref{sec:neural} described a code whose active units carry
a category structure; the natural next question from the coding side is
whether that structure \emph{factors} --- whether the code decomposes into
independent groups of units, each carrying its own small DAG, so that an
item's description is a tuple rather than a set. That is a factorial code in
Barlow's sense \citep{barlow1989}, a product lattice in order-theoretic terms,
and a product of pattern structures in \fca{} terms.

\odag{} already has a degenerate version: its dimensions are exactly such
independent factors (weight, time and location vary independently, and a value
in one says nothing about the others). Generalising from a handful of declared
dimensions to \emph{learned} factors would connect the MDL programme to the
redundancy-reduction tradition it came from, and would give the surface layer
a principled account of when two categories are ``about the same thing''.
The machine-learning neighbours to read first are the layer-wise
total-correlation approach of \citet{versteeg2014} and the unbounded-depth
Bayesian prior over binary latent DAGs of \citet{adams2010}.

\paragraph{D15. Continuous relaxations as a recall layer.} Order embeddings
\citep{vendrov2016} and hyperbolic embeddings \citep{nickel2017} learn vectors
in which subsumption is coordinate dominance or radial position. They are the
approximate, differentiable cousins of the ancestor-set code, and they fail in
exactly the way the exact structure does not: approximately. That makes them a
natural \emph{proposal} layer over an exact \emph{disposal} layer, which is
precisely the division of labour Section~\ref{sec:agents} recommends for
retrieval. Worth noting that the tree/DAG obstruction reappears there too ---
trees embed in hyperbolic space with arbitrarily low distortion, DAGs only
approximately.

\paragraph{D16. Closing the mind--brain loop.} The most speculative and, for
this project's origins, the most fitting. Figure~\ref{fig:mind-brain} has two
directions and the literature has mostly travelled one of them: from recorded
codes to concept lattices \citep{endres2008,endres2010}. The other direction
is now buildable: take an ontology, embed it as an ancestor-set code
(Proposition~\ref{prop:roundtrip} says nothing is lost), and ask whether a
population code with that structure predicts responses better than one
without. \odag{} would be supplying the ontology in a form where the embedding
is exact and reproducible --- which is to say, canonical --- and where two
laboratories could confirm they used the identical ontology by comparing 32
bytes.

\subsection{The map}

\begin{table}[h]\centering\small
\caption{All directions, with the trigger that would justify starting.}
\label{tab:directions}
\vspace{1mm}
\begin{tabularx}{\textwidth}{@{}p{0.7cm} p{4.5cm} p{2.2cm} X@{}}
\toprule
& \textbf{Direction} & \textbf{Risk to invariants} & \textbf{Trigger}\\
\midrule
D1 & Pattern-structure framing & none & now (free)\\
D2 & Ordinal kind & none & someone files ranks\\
D3 & Reduced labelling & none & now\\
D4 & Implication lint & none (derived) & users mis-filing\\
D5 & Cone bitmaps & none (derived) & queries hot\\
\midrule
D6 & Exploration with an LLM expert & none (claims) & now; the best fit with
     agents\\
D7 & Workload materialisation & soundness caveat & a real query log\\
D8 & MDL concepts + promotion & none (derived) & store too big to see\\
D9 & Merkle cone commitments & none (derived) & on-chain subsumption needed\\
D10 & Thin browser client & none & a published store worth reading\\
\midrule
D11 & Concurrent exploration & unclear & multi-writer exploration\\
D12 & Relations & \textbf{breaks canonicality} & mass reification observed\\
D13 & Zero-knowledge subsumption & none & a private counterparty\\
D14 & Products of DAGs & none (derived) & research\\
D15 & Embeddings as recall & none (outside) & vague-input queries\\
D16 & Ontology $\to$ neural code & none & a collaborator with data\\
\bottomrule
\end{tabularx}
\end{table}

Only one row in that table says ``breaks canonicality'', and it is the one
everybody asks for first. That is the paper's practical warning, and
Section~\ref{sec:conclusion} makes it the closing point.
